\chapter{Appendix: capability detection cheatsheet}
\label{ch:capref}

\begin{aside}
The escapes \maya{} sends at startup to detect what the
terminal supports.
\end{aside}

\section{Truecolour}

\code{COLORTERM=truecolor} (env). Send a truecolour SGR and
query position; some terminals echo. \maya{} trusts env.

\section{256-colour}

\code{TERM} contains \code{256color}. Or query
\code{ESC [ ? 1 c}; a modern terminal reports.

\section{Image protocols}

\begin{itemize}[leftmargin=*]
  \item Kitty: query \code{ESC \_ G i=1,s=1,v=1,a=q ; ESC \textbackslash}.
  \item Sixel: device attributes \code{ESC [ c} response
    includes 4.
  \item iTerm2: \code{TERM\_PROGRAM=iTerm.app} (env).
\end{itemize}

\section{OSC-8 hyperlinks}

No detection escape. Send the sequence; terminals that
don't support it ignore. \maya{} sends unconditionally.

\section{CSI-u}

Send \code{ESC [ > 1 u}; query with \code{ESC [ ? u}.
Modern terminals reply.

\section{DEC 2026 synchronised output}

Query mode 2026 with \code{ESC [ ? 2026 \$ p}. Reply
\code{ESC [ ? 2026 ; 0/1/2 \$ y}.

\section{Mouse SGR}

Send \code{ESC [ ? 1006 h}; no reply. Try mouse events;
if SGR comes back, supported.

\section{Bracketed paste}

\code{ESC [ ? 2004 h}; no reply. Universal on modern
terminals.

\section{Pixel size}

\code{ESC [ 16 t} reports cell-pixel size as
\code{ESC [ 6 ; H ; W t}. Used for inline image sizing.

\section{Recap}

\begin{recap}
\begin{itemize}[leftmargin=*]
  \item Env vars first; queries second.
  \item Some features are query-and-listen, others are
    fire-and-hope.
  \item Cache results for the session.
\end{itemize}
\end{recap}
